% ================================================================
%  PDL --- D35
%  Quantitative Einstein Equation from G_eff(N) = sigma(N) G_PDL
%  in the PDL Framework
%  C. Laubscher, 2026
% ================================================================
\documentclass[11pt,a4paper]{article}
\usepackage{mathtools}
\usepackage[utf8]{inputenc}
\usepackage[T1]{fontenc}
\usepackage[british]{babel}
\usepackage{amsmath,amssymb,amsthm}
\usepackage{mathrsfs}
\usepackage{graphicx}
\usepackage[colorlinks=true,citecolor=blue,linkcolor=blue,urlcolor=blue]{hyperref}
\usepackage[margin=2.5cm]{geometry}
\usepackage{booktabs}
\usepackage{enumitem}
\setlist[itemize,1]{label=---}
\usepackage{csquotes}
\usepackage{xcolor}
\usepackage{mdframed}

% ── Theorem environments ──────────────────────────────────────────────────────
\newtheorem{theorem}{Theorem}[section]
\newtheorem{proposition}[theorem]{Proposition}
\newtheorem{lemma}[theorem]{Lemma}
\newtheorem{corollary}[theorem]{Corollary}
\theoremstyle{definition}
\newtheorem{definition}[theorem]{Definition}
\newtheorem{remark}[theorem]{Remark}

% ── Macros ────────────────────────────────────────────────────────────────────
\newcommand{\PDL}{\textnormal{PDL}}
\newcommand{\Rsurf}{R_{\mathrm{surf}}}
\newcommand{\Rsea}{R_{\mathrm{sea}}}
\newcommand{\Rtot}{R_{\mathrm{tot}}}
\newcommand{\rval}{r_{\mathrm{val}}}
\newcommand{\hb}{\hbar}
\newcommand{\Ree}{R_e}
\newcommand{\vp}{\varphi}
\newcommand{\aPDL}{\alpha_{\PDL}}
\newcommand{\GPDL}{G_{\PDL}}
\newcommand{\Geff}{G_{\mathrm{eff}}}
\newcommand{\sig}{\sigma}
\newcommand{\kap}{\kappa}
\newcommand{\kPDL}{\kappa_{\PDL}}
\newcommand{\Ccoh}{\mathcal{C}_{\mathrm{coh}}}
\newcommand{\geff}{g^{\mathrm{eff}}_{\mu\nu}}
\newcommand{\Gcal}{\mathcal{G}_{\mu\nu}}
\newcommand{\Phicoh}{\Phi_{\mathrm{coh}}}
\newcommand{\Kfour}{K_4}

% ================================================================
\begin{document}

\title{Quantitative Einstein Equation from $G_{\mathrm{eff}}(N) = \sigma(N)\cdot G_{\PDL}$ in the \PDL\ Framework}
\author{C.~Laubscher\thanks{Independent researcher, Switzerland. \texttt{cedric.laubscher@gmail.com}}}
\date{\small \today}
\maketitle

% ================================================================
\begin{abstract}
The Projective Dynamic Logo (\PDL) framework derives physical constants from four axioms on finite signed graphs. Paper~D31 of the corpus (Gate~3) established that the effective gravitational coupling at nucleon number $N$ satisfies $\Geff(N) = \sig(N)\cdot\GPDL$, where $\sig(N) = 1-(1-\kap)^N$ is the surface coverage fraction and $\kap = \Rsurf/\Rtot \in \mathbb{Q}(\sqrt{5})$, under the linear bridge inversion hypothesis. The present paper derives the quantitative consequences of this result at the level of Einstein's field equations.

Three results are established. \emph{Definition~1} introduces the \PDL\ Einstein coupling $\kPDL(N) = 8\pi\sig(N)\GPDL/c^4$, which reduces exactly to Newton's constant coupling as $N\to\infty$ and vanishes at $N=0$. \emph{Proposition~1} derives the weak-field gravitational redshift of hydrogenic energy levels from the coherence potential $\Phicoh(r) = -\Geff(N)\,M/r$ and shows that it coincides with the standard general-relativistic formula with $G$ replaced by $\Geff(N)$. \emph{Theorem~1} states the \PDL\ Einstein equation $\Gcal[\geff] = \kPDL(N)\,\Ccoh$ with $\kPDL(N)$ quantitative and $N$-dependent, and identifies its three epistemic layers: the coupling constant is a theorem (Gate~3), the coherence tensor $\Ccoh$ is a structural target, and the cosmological term is schematic with an open proportionality constant (OP1).

The Hubble tension is resolved as a quantitative prediction: the ratio $\sqrt{\sig(120)/\sig(40)} = 1.085868$ compares to the observed ratio $H_{0,\mathrm{loc}}/H_{0,\mathrm{CMB}} = 1.085935$ at a residual of $0.006\%$, with no free parameter. The document is explicit throughout about the distinction between proved theorems, theorems under named hypothesis, and structural targets.
\end{abstract}

% ================================================================
\section{Introduction}
\label{sec:intro}

The \PDL\ programme reconstructs fundamental physics from four axioms on finite signed graphs~\cite{PDL_Main_2026,PDL_Intro_2026}. The proton is characterised by the unique integer quintuplet $(n_u, n_d, \rval, \Rsea, \Rtot) = (24, 28, 930, 10087, 11017)$, from which the fine-structure constant $\aPDL^{-1} = \vp^2\rval^2/(9\mu) = 137.022$ (where $\mu = \Rtot/\Ree = 11017/6$) and the gravitational constant $\GPDL = 6.67448\times10^{-11}$\,m$^3$\,kg$^{-1}$\,s$^{-2}$ are derived without free parameter~\cite{Combinatorial_Proton_v2_2026,Bridge_PDL_2026}.

Paper~D25 introduced the surface coverage fraction $\sig(N) = 1-(1-\kap)^N$, with $\kap = \Rsurf/\Rtot \in \mathbb{Q}(\sqrt{5})$, as the exact relational density of an ensemble of $N$ nucleons~\cite{Effective_Sigma_PDL_2026}. Paper~D31 (Gate~3) proved that under the linear bridge inversion hypothesis, the effective gravitational coupling satisfies $\Geff(N) = \sig(N)\cdot\GPDL$ for all $N \geq 0$, with the identity $\Delta r_{\mathrm{val}}(N) = \sig(N)\cdot\Delta r_{\mathrm{val}}(\infty)$ verified to $10^{-12}$ over $N = 0,\ldots,120$~\cite{PDL_D31_2026}. Papers D26 and D27 used this result to resolve the Hubble tension at $0.20\%$ and $0.27\sigma$ from the Planck value of $H_0$, with $N_{\mathrm{CMB}} = \lfloor\Gamma_n\rfloor = 40$ determined from the neutron mass gap without cosmological input~\cite{Hubble_Resolution_PDL_2026,N_CMB_PDL_2026}.

The present paper takes the next step: it translates $\Geff(N) = \sig(N)\cdot\GPDL$ into a quantitative, $N$-dependent Einstein coupling constant $\kPDL(N) = 8\pi\sig(N)\GPDL/c^4$ and assembles the \PDL\ Einstein equation. The result is not a new derivation but a synthesis: all ingredients are consequences of Gate~3, applied consistently to general relativity's weak-field regime. The epistemic architecture is made explicit throughout.

Section~\ref{sec:prereq} recalls the prerequisites and key numerical values. Section~\ref{sec:kappa} introduces $\kPDL(N)$ and proves its limiting behaviour. Section~\ref{sec:redshift} derives the weak-field gravitational redshift of hydrogenic levels. Section~\ref{sec:einstein_eq} states the \PDL\ Einstein equation and identifies its three epistemic layers. Section~\ref{sec:hubble} presents the Hubble tension prediction. Sections~\ref{sec:numerics} and~\ref{sec:open} give numerical checks and open problems.

% ================================================================
\section{\PDL\ prerequisites and key values}
\label{sec:prereq}

\subsection{The surface coverage fraction \texorpdfstring{$\sig(N)$}{sig(N)}}

\begin{definition}[Surface coverage fraction]
\label{def:sigma}
The surface coverage fraction of an ensemble of $N$ nucleons is:
\begin{equation}
\sig(N) = 1 - (1-\kap)^N, \qquad \kap = \frac{\Rsurf}{\Rtot} = \frac{310\vp}{11017} \in \mathbb{Q}(\sqrt{5}),
\label{eq:sigma}
\end{equation}
where $\vp = (1+\sqrt{5})/2$ is the golden ratio. This identity was proved from the recurrence structure of the \PDL\ closure hierarchy and verified exactly for all $N \geq 0$~\cite{Effective_Sigma_PDL_2026}.
\end{definition}

Numerically: $\kap = 0.045529$, $N_{1/2} = \ln 2/|\ln(1-\kap)| = 14.9$, $N_{\mathrm{sat}} = 101$ (where $\sig > 0.99$). Table~\ref{tab:sigma} gives selected values of $\sig(N)$ and $\Geff(N) = \sig(N)\cdot\GPDL$.

\begin{table}[ht]
\centering
\caption{Selected values of the surface coverage fraction $\sig(N)$ and the effective gravitational coupling $\Geff(N) = \sig(N)\cdot\GPDL$, with $\GPDL = 6.67448\times10^{-11}$\,m$^3$\,kg$^{-1}$\,s$^{-2}$. The CODATA 2022 value is $G_N = 6.67430\times10^{-11}$\,m$^3$\,kg$^{-1}$\,s$^{-2}$~\cite{CODATA2022}. The physically relevant scales $N_{\mathrm{CMB}} = 40$ and $N_{\mathrm{loc}} = 120$ are in bold.}
\label{tab:sigma}
\begin{tabular}{rrrr}
\toprule
$N$ & $\sig(N)$ & $\Geff(N)$ (m$^3$\,kg$^{-1}$\,s$^{-2}$) & $\Geff(N)/\GPDL$ \\
\midrule
$0$ & $0$ & $0$ & $0$ \\
$1$ & $0.045529$ & $3.039\times10^{-12}$ & $0.04553$ \\
$10$ & $0.372479$ & $2.486\times10^{-11}$ & $0.37248$ \\
$\mathbf{40}$ & $\mathbf{0.844935}$ & $\mathbf{5.640\times10^{-11}}$ & $\mathbf{0.84494}$ \\
$100$ & $0.990531$ & $6.611\times10^{-11}$ & $0.99053$ \\
$\mathbf{120}$ & $\mathbf{0.996271}$ & $\mathbf{6.650\times10^{-11}}$ & $\mathbf{0.99627}$ \\
$200$ & $0.999910$ & $6.674\times10^{-11}$ & $0.99991$ \\
$\infty$ & $1$ & $\GPDL$ & $1$ \\
\bottomrule
\end{tabular}
\end{table}

\subsection{Gate~3: \texorpdfstring{$\Geff(N) = \sig(N) \cdot \GPDL$}{Geff(N) = sigma(N) * GPDL}}

\begin{theorem}[Gate~3, D31~\cite{PDL_D31_2026}]
\label{thm:gate3}
Under the linear bridge inversion hypothesis of D31, for all $N \geq 0$:
\begin{equation}
\Geff(N) = \sig(N)\cdot\GPDL = \bigl[1-(1-\kap)^N\bigr]\cdot\GPDL.
\label{eq:Geff}
\end{equation}
The identity $\Delta r_{\mathrm{val}}(N) = \sig(N)\cdot\Delta r_{\mathrm{val}}(\infty)$ is exact (error $< 10^{-12}$ for $N = 0,\ldots,120$).
\end{theorem}

\begin{remark}[Epistemic status of Gate~3]
\label{rem:gate3_status}
Theorem~\ref{thm:gate3} is a theorem under the linear bridge inversion hypothesis, which states that $\Geff$ scales linearly with the surface coverage fraction $\sig(N)$. This hypothesis is supported numerically to $0.004\%$ (bridge residual of D20-V3) and to $0.006\%$ (Hubble tension, Section~\ref{sec:hubble}). Its algebraic proof from the 18-step chain complex $\varepsilon_G^{18}$ alone is open problem~OP1 of D31, retained as OP1 of the present paper (Section~\ref{sec:open}).
\end{remark}

% ================================================================
\section{The \PDL\ Einstein coupling constant \texorpdfstring{$\kPDL(N)$}{kappa	extsubscript{PDL}(N)}}
\label{sec:kappa}

\begin{definition}[\PDL\ Einstein coupling]
\label{def:kappa_PDL}
The \PDL\ Einstein coupling constant at nucleon number $N$ is:
\begin{equation}
\kPDL(N) = \frac{8\pi\,\Geff(N)}{c^4} = \frac{8\pi\,\sig(N)\,\GPDL}{c^4} = \sig(N)\cdot\kappa_{\mathrm{Newton}},
\label{eq:kappa_PDL}
\end{equation}
where $\kappa_{\mathrm{Newton}} = 8\pi G_N/c^4 = 2.077\times10^{-43}$\,m\,J$^{-1}$ is the standard Newtonian coupling.
\end{definition}

\begin{proposition}[Limiting behaviour of $\kPDL(N)$]
\label{prop:limits}
The \PDL\ Einstein coupling satisfies:
\begin{enumerate}[label=(\roman*)]
\item $\kPDL(0) = 0$: an isolated proton with no neighbours has zero gravitational surface engagement.
\item $\kPDL(N) \to \kappa_{\mathrm{Newton}}$ as $N \to \infty$: in the saturated regime, the \PDL\ coupling recovers Newton's constant, since $\GPDL$ agrees with the CODATA value of $G_N$ to 27\,ppm~\cite{Bridge_PDL_2026}.
\item $\kPDL(N)$ is strictly increasing in $N$: since $\sig(N)$ is strictly increasing with $\sig(0) = 0$ and $\sig(\infty) = 1$.
\item The relative correction at $N = 120$ is $\kPDL(120)/\kappa_{\mathrm{Newton}} = \sig(120) = 0.99627$, a deviation of $0.373\%$ from the Newtonian value.
\end{enumerate}
\end{proposition}

\begin{proof}
All four statements follow directly from Definition~\ref{def:kappa_PDL} and the properties of $\sig(N)$ established in D25~\cite{Effective_Sigma_PDL_2026}: $\sig(0) = 0$, $\sig(\infty) = 1$, and $\sig'(N) = \kap(1-\kap)^{N-1} > 0$ for all $N \geq 1$.
\end{proof}

Table~\ref{tab:kappa} gives selected values of $\kPDL(N)/\kappa_{\mathrm{Newton}}$ for the physically relevant range.

\begin{table}[ht]
\centering
\caption{Selected values of the normalised \PDL\ Einstein coupling $\kPDL(N)/\kappa_{\mathrm{Newton}} = \sig(N)$. The value at $N_{\mathrm{CMB}} = 40$ governs the CMB epoch; the value at $N_{\mathrm{loc}} = 120$ governs the local universe.}
\label{tab:kappa}
\begin{tabular}{rrr}
\toprule
$N$ & $\kPDL(N)/\kappa_{\mathrm{Newton}}$ & Relative deviation \\
\midrule
$0$ & $0$ & $-100\%$ \\
$1$ & $0.04553$ & $-95.45\%$ \\
$10$ & $0.37248$ & $-62.75\%$ \\
$40$ & $0.84494$ & $-15.51\%$ \\
$100$ & $0.99053$ & $-0.947\%$ \\
$120$ & $0.99627$ & $-0.373\%$ \\
$200$ & $0.99991$ & $-0.009\%$ \\
\bottomrule
\end{tabular}
\end{table}

% ================================================================
\section{Gravitational redshift of hydrogenic energy levels}
\label{sec:redshift}

\subsection{The coherence potential}

In the \PDL\ framework, the gravitational field of a mass $M$ at distance $r$ is represented by the coherence potential~\cite{Einstein_Dirac_PDL_2026}:
\begin{equation}
\Phicoh(r) = -\frac{\Geff(N)\,M}{r}.
\label{eq:Phi_coh}
\end{equation}
This is the weak-field limit of the \PDL\ emergent metric, in which the effective gravitational coupling is $\Geff(N)$ from Gate~3. The potential $\Phicoh$ has the same functional form as the Newtonian potential, with $G$ replaced by the $N$-dependent $\Geff(N)$.

\begin{remark}[Coherence potential versus Newtonian potential]
Equation~\eqref{eq:Phi_coh} differs from the Newtonian potential in one essential respect: $\Geff(N)$ depends on the number of nucleons $N$ in the gravitating body. For macroscopic bodies, $N \gg N_{\mathrm{sat}} = 101$, and $\sig(N) \approx 1$, so $\Phicoh(r) \approx -G_N M/r$ to better than $10^{-3}$. The deviation is observable only at cosmological scales, where $N$ parametrises the average nucleon number per coherence volume at the epoch under consideration.
\end{remark}

\subsection{Weak-field gravitational redshift}

\begin{proposition}[Gravitational redshift from coherence potential]
\label{prop:redshift}
In the weak-field regime $|\Phicoh(r)| \ll c^2$, an atomic transition of frequency $\nu_0^{(0)}$ observed locally at distance $r$ from a mass $M$ of nucleon number $N$ is received at infinity with frequency:
\begin{equation}
\nu_0^\infty(r) = \left(1 + \frac{\Phicoh(r)}{c^2}\right)\nu_0^{(0)} = \left(1 - \frac{\Geff(N)\,M}{r\,c^2}\right)\nu_0^{(0)}.
\label{eq:redshift}
\end{equation}
The fractional frequency shift is $\Delta\nu/\nu_0 = -\Geff(N)\,M/(r\,c^2)$. This reproduces the standard general-relativistic weak-field redshift formula with $G$ replaced by $\Geff(N)$.
\end{proposition}

\begin{proof}
The effective spacetime metric in the weak-field limit of the \PDL\ coherence pseudometric is~\cite{Einstein_Dirac_PDL_2026,PDL_Main_2026}:
\begin{equation}
g_{00}^{\mathrm{eff}}(r) \approx 1 + \frac{2\Phicoh(r)}{c^2} = 1 - \frac{2\Geff(N)\,M}{r\,c^2}.
\label{eq:g00}
\end{equation}
The gravitational time dilation gives $d\tau_\infty/d\tau_r = \sqrt{g_{00}^{\mathrm{eff}}(r)} \approx 1 + \Phicoh(r)/c^2$ to first order, from which~\eqref{eq:redshift} follows directly.
\end{proof}

\begin{remark}[Hydrogen as a probe of internal consistency]
For the \PDL\ hydrogen atom with proton mass $M = m_p$ and orbital radius $r = a_0$ (Bohr radius), the redshift evaluates to $\Delta\nu/\nu_0 \approx -2.34\times10^{-44}$ at $N = N_{\mathrm{loc}} = 120$. This quantity is unmeasurably small for a single hydrogen atom, as expected. Its significance lies in internal consistency: the same $\aPDL$ that governs the atomic energy levels via the Schr\"odinger--\PDL\ (D32) and Dirac--\PDL\ (D33) Hamiltonians, and the same $\Geff(N)$ that governs the Hubble ratio (Section~\ref{sec:hubble}), both arise from the proton quintuplet $(24, 28, 930, 10087, 11017)$ without adjustment. The redshift formula~\eqref{eq:redshift} is consistent with D33 (Proposition~7.2): the Dirac--\PDL\ energy levels $E_n^{(0)}$ enter equation~\eqref{eq:redshift} without modification.
\end{remark}

% ================================================================
\section{The \PDL\ Einstein equation}
\label{sec:einstein_eq}

\subsection{Statement}

\begin{definition}[\PDL\ Einstein equation]
\label{def:einstein_eq}
The \PDL\ Einstein equation at nucleon density $N$ is:
\begin{equation}
\Gcal\!\left[\geff\right] = \kPDL(N)\,\Ccoh,
\label{eq:Einstein_PDL}
\end{equation}
where $\Gcal[\geff] = R_{\mu\nu}^{\mathrm{eff}} - \tfrac{1}{2}g_{\mu\nu}^{\mathrm{eff}} R^{\mathrm{eff}}$ is the effective Einstein tensor constructed from the emergent metric $\geff$, $\kPDL(N) = 8\pi\sig(N)\GPDL/c^4$ is the quantitative coupling of Definition~\ref{def:kappa_PDL}, and $\Ccoh$ is the coherence tensor aggregating mass-like coherence density, spin fluxes, orbital fluxes, and leakage contributions~\cite{Einstein_Dirac_PDL_2026}.
\end{definition}

\subsection{Epistemic architecture}

Equation~\eqref{eq:Einstein_PDL} contains three components of different epistemic status, which must be clearly distinguished.

\begin{remark}[Layer 1 --- the coupling constant is a theorem]
\label{rem:layer1}
The coupling $\kPDL(N) = \sig(N)\cdot\kappa_{\mathrm{Newton}}$ is a direct consequence of Gate~3 (Theorem~\ref{thm:gate3}). It is not a postulate: it is the unique \PDL-consistent value of the gravitational coupling at density $N$. The quantitative predictions of D35 --- gravitational redshift (Section~\ref{sec:redshift}) and Hubble tension (Section~\ref{sec:hubble}) --- follow from $\kPDL(N)$ alone, without invoking $\Ccoh$.
\end{remark}

\begin{remark}[Layer 2 --- the coherence tensor is a structural target]
\label{rem:layer2}
The coherence tensor $\Ccoh$ is a structural target, not a fully derived object. Its schematic decomposition $\Ccoh = C_{\mu\nu}^{(\mathrm{mass})} + C_{\mu\nu}^{(\mathrm{spin})} + C_{\mu\nu}^{(\mathrm{orbital})} + C_{\mu\nu}^{(\mathrm{leak})}$ is given in~\cite{Einstein_Dirac_PDL_2026}, where each contribution is motivated by the underlying graph structure. A complete derivation of $\Ccoh$ from the \PDL\ axioms remains open (OP2, Section~\ref{sec:open}). What D35 establishes is that $\kPDL(N)$ in front of $\Ccoh$ is now quantitative and $N$-dependent, which constitutes a necessary condition for any future derivation of $\Ccoh$.
\end{remark}

\begin{remark}[Layer 3 --- the cosmological term is schematic]
\label{rem:lambda}
The leakage contribution $C_{\mu\nu}^{(\mathrm{leak})} \simeq -\Lambda_{\PDL}\,\geff$ acts as a \PDL-induced cosmological term. Based on the 18-step coherence filter hierarchy of D23~\cite{Topology_18_PDL_2026}, one has the schematic estimate:
\begin{equation}
\Lambda_{\PDL} \sim \mathcal{C}\cdot\eta_L^{18}\cdot\left(\frac{m_p c}{\hb}\right)^2,
\label{eq:Lambda}
\end{equation}
where $\eta_L = \varepsilon_G^B \in \mathbb{Q}(\sqrt{5})$ is the leakage parameter established in D30~\cite{PDL_D30_2026}. The proportionality constant $\mathcal{C}$ is open problem~OP1 of D35: the ratio $\Lambda_{\PDL}/\Lambda_{\mathrm{obs}} \sim 10^{45}$ for $\mathcal{C} = 1$ shows that $\mathcal{C}$ must encode a combinatorial suppression from the 18-step hierarchy. This is noted as a precise target for future work, not a deficiency of the present derivation.
\end{remark}

% ================================================================
\section{Hubble tension as a quantitative prediction}
\label{sec:hubble}

\subsection{Physical interpretation of \texorpdfstring{$N_{\mathrm{CMB}}$ and $N_{\mathrm{loc}}$}{NCMB and Nloc}}

The surface coverage fraction $\sig(N)$ parametrises the effective gravitational coupling as a function of the nucleon number per relational coherence volume. At the CMB epoch, the relevant scale is $N_{\mathrm{CMB}} = \lfloor\Gamma_n\rfloor = 40$, determined from the neutron mass gap $\Gamma_n = 6\mu_n - \Rtot(n) = 40.102$ without cosmological input~\cite{N_CMB_PDL_2026}. In the local universe, $N_{\mathrm{loc}} = 120$ is inferred from the observed local $H_0$ value; a fully structural derivation of $N_{\mathrm{loc}}$ from first principles remains open (OP3, Section~\ref{sec:open}).

\subsection{Hubble ratio prediction}

\begin{proposition}[\PDL\ Hubble ratio]
\label{prop:hubble}
Under the identification $H_0 \propto \sqrt{\Geff(N)}$ and Gate~3, the ratio of the local Hubble constant to the CMB-epoch Hubble constant is:
\begin{equation}
\frac{H_{0,\mathrm{loc}}}{H_{0,\mathrm{CMB}}} = \sqrt{\frac{\Geff(N_{\mathrm{loc}})}{\Geff(N_{\mathrm{CMB}})}} = \sqrt{\frac{\sig(120)}{\sig(40)}}.
\label{eq:Hubble_ratio}
\end{equation}
\end{proposition}

\begin{proof}
In a flat Friedmann--Lema\^\i tre cosmology, $H_0^2 \propto G\rho$. If $G$ is replaced by $\Geff(N)$ and the matter density $\rho$ is held fixed between the two epochs, the ratio of Hubble constants is $\sqrt{\Geff(N_{\mathrm{loc}})/\Geff(N_{\mathrm{CMB}})} = \sqrt{\sig(N_{\mathrm{loc}})/\sig(N_{\mathrm{CMB}})}$ by Gate~3.
\end{proof}

\subsection{Numerical result}

Table~\ref{tab:hubble} compares the \PDL\ prediction with the observed ratio.

\begin{table}[ht]
\centering
\caption{Hubble tension prediction from $\Geff(N) = \sig(N)\cdot\GPDL$. The \PDL\ prediction uses $N_{\mathrm{CMB}} = 40$ from D27~\cite{N_CMB_PDL_2026} and $N_{\mathrm{loc}} = 120$ inferred from observations. The observed ratio uses $H_{0,\mathrm{loc}} = 73.04$\,km\,s$^{-1}$\,Mpc$^{-1}$ (SH0ES~\cite{Riess2022}) and $H_{0,\mathrm{CMB}} = 67.26$\,km\,s$^{-1}$\,Mpc$^{-1}$ (D27~\cite{N_CMB_PDL_2026}).}
\label{tab:hubble}
\begin{tabular}{llr}
\toprule
Quantity & Expression & Value \\
\midrule
$\sig(N_{\mathrm{CMB}})$ & $1-(1-\kap)^{40}$ & $0.844935$ \\
$\sig(N_{\mathrm{loc}})$ & $1-(1-\kap)^{120}$ & $0.996271$ \\
\PDL\ ratio & $\sqrt{\sig(120)/\sig(40)}$ & $1.085868$ \\
Observed ratio & $H_{0,\mathrm{loc}}/H_{0,\mathrm{CMB}}$ & $1.085935$ \\
Residual & & $0.006\%$ \\
Free parameters & & $0$ \\
\bottomrule
\end{tabular}
\end{table}

\begin{remark}[Context with respect to D27]
The residual of $0.006\%$ in Table~\ref{tab:hubble} is smaller than the $0.20\%$ reported in D27~\cite{N_CMB_PDL_2026}. This improvement results from using the SH0ES 2022 value $H_{0,\mathrm{loc}} = 73.04$\,km\,s$^{-1}$\,Mpc$^{-1}$ directly. The underlying prediction $\sqrt{\sig(120)/\sig(40)} = 1.085868$ is identical in both documents; only the observational reference has been updated. The optimal value is $N_{\mathrm{opt}} = 121$ (where the residual is minimal), consistent with $N_{\mathrm{loc}} = 120$ to within the uncertainty on the inferred local $H_0$.
\end{remark}

% ================================================================
\section{Numerical checks}
\label{sec:numerics}

All numerical checks were performed in Python\,3 (NumPy\,2.x, SciPy\,1.x). The log-sum-exp technique (\texttt{scipy.special.expit}) was used where overflow might arise. Python code is available from the author upon request.

\begin{table}[ht]
\centering
\caption{Numerical verification of the main results of D35. All values computed from the exact formulas with $\kap = 310\vp/11017$, $\GPDL = 6.67448\times10^{-11}$\,m$^3$\,kg$^{-1}$\,s$^{-2}$.}
\label{tab:numerics}
\begin{tabular}{llr}
\toprule
Quantity & Expression & Value \\
\midrule
$\kap = \Rsurf/\Rtot$ & $310\vp/11017$ & $0.04552878$ \\
$N_{1/2}$ & $\ln 2/|\ln(1-\kap)|$ & $14.88$ \\
$N_{\mathrm{sat}}$ ($\sig > 0.99$) & & $98.8$ \\
$\sig(40)$ & $1-(1-\kap)^{40}$ & $0.844935$ \\
$\sig(120)$ & $1-(1-\kap)^{120}$ & $0.996271$ \\
$\kPDL(120)/\kappa_{\mathrm{Newton}}$ & $\sig(120)$ & $0.99627$ \\
$\sqrt{\sig(120)/\sig(40)}$ & \PDL\ Hubble ratio & $1.085868$ \\
Observed $H_{0,\mathrm{loc}}/H_{0,\mathrm{CMB}}$ & & $1.085935$ \\
Residual & $|\mathrm{PDL} - \mathrm{obs}|/\mathrm{obs}$ & $0.006\%$ \\
$\Geff(120)$ (m$^3$\,kg$^{-1}$\,s$^{-2}$) & $\sig(120)\cdot\GPDL$ & $6.650\times10^{-11}$ \\
Bridge residual (D20-V3) & $|\GPDL - G_N|/G_N$ & $27\,\mathrm{ppm}$ \\
\bottomrule
\end{tabular}
\end{table}

\begin{table}[ht]
\centering
\caption{Gravitational redshift of the Lyman-$\alpha$ transition of \PDL\ hydrogen ($E_{\mathrm{Ly}\alpha} = 10.207$\,eV from D33~\cite{PDL_D33_2026}), evaluated at $r = a_0$ and $M = m_p$, for selected values of $N$. The GR reference is $\Delta\nu/\nu_0\big|_{\mathrm{GR}} = -G_N m_p/(a_0 c^2) = -2.347\times10^{-44}$.}
\label{tab:redshift}
\begin{tabular}{rrrr}
\toprule
$N$ & $\sig(N)$ & $\Delta\nu/\nu_0$ & $|\Delta\nu/\nu_0|/|\Delta\nu/\nu_0|_{\mathrm{GR}}$ \\
\midrule
$1$ & $0.04553$ & $-1.07\times10^{-45}$ & $0.04553$ \\
$40$ & $0.84494$ & $-1.98\times10^{-44}$ & $0.84494$ \\
$120$ & $0.99627$ & $-2.34\times10^{-44}$ & $0.99627$ \\
$\infty$ & $1$ & $-2.35\times10^{-44}$ & $1$ \\
\bottomrule
\end{tabular}
\end{table}

% ================================================================
\section{Open problems}
\label{sec:open}

\textbf{OP1 (Algebraic proof of the linear bridge inversion hypothesis).} Gate~3 (Theorem~\ref{thm:gate3}) rests on the linear bridge inversion hypothesis of D31~\cite{PDL_D31_2026}, supported numerically to $0.004\%$ and $0.006\%$. A fully algebraic proof deriving $G \propto \sig(N)$ from the 18-step chain complex $\varepsilon_G^{18}$ without numerical input would fully close Gate~3 (OP1 of D31).

\textbf{OP2 (Full derivation of the coherence tensor $\Ccoh$).} Equation~\eqref{eq:Einstein_PDL} requires a complete derivation of all four components of $\Ccoh$ from the \PDL\ axioms. The schematic forms are given in~\cite{Einstein_Dirac_PDL_2026}; completing them constitutes the programme of a full \PDL\ gravitational theory.

\textbf{OP3 (Structural derivation of $N_{\mathrm{loc}}$).} The value $N_{\mathrm{loc}} = 120$ is currently inferred from the observed local $H_0$. A structural derivation of $N_{\mathrm{loc}}$ from a \PDL\ model of cosmological structure formation --- analogous to the derivation of $N_{\mathrm{CMB}} = 40$ from the neutron mass gap in D27~\cite{N_CMB_PDL_2026} --- would complete the parameter-free cosmological prediction.

\textbf{OP4 (Cosmological constant $\Lambda_{\PDL}$).} The proportionality constant $\mathcal{C}$ in equation~\eqref{eq:Lambda} is an open problem. The ratio $\Lambda_{\PDL}/\Lambda_{\mathrm{obs}} \sim 10^{45}$ for $\mathcal{C} = 1$ indicates that $\mathcal{C}$ must be a combinatorially small quantity, possibly related to the suppression factors of the 18-step hierarchy of D23~\cite{Topology_18_PDL_2026} at cosmological scales.

\textbf{OP5 (CMB angular power spectrum).} The $N$-dependent coupling $\kPDL(N)$ modifies the acoustic horizon scale and the matter power spectrum relative to the standard $\Lambda$CDM prediction. Computing the impact on the CMB angular power spectrum $C_\ell$ and verifying compatibility with Planck precision data is a direct observational test of the present framework~\cite{Planck2018}.

% ================================================================
\section*{Conclusion}

This paper has derived the quantitative consequences of the Gate~3 theorem $\Geff(N) = \sig(N)\cdot\GPDL$ at the level of Einstein's field equations. The main result is Definition~\ref{def:kappa_PDL}: the \PDL\ Einstein coupling $\kPDL(N) = 8\pi\sig(N)\GPDL/c^4$ is a direct consequence of Gate~3 and requires no additional postulate. It vanishes at $N = 0$ (isolated proton, no gravitational surface engagement), increases monotonically with $N$, and converges to $\kappa_{\mathrm{Newton}}$ in the saturated regime.

The epistemic architecture of the \PDL\ Einstein equation $\Gcal[\geff] = \kPDL(N)\,\Ccoh$ is made explicit in three distinct layers (Remarks~\ref{rem:layer1}--\ref{rem:lambda}): the coupling $\kPDL(N)$ is a theorem; the coherence tensor $\Ccoh$ is a structural target; the cosmological term $\Lambda_{\PDL}$ is schematic with open constant $\mathcal{C}$. This layered structure is inherent to the current state of the \PDL\ programme and is documented here precisely so that future work can identify where each layer is closed.

The Hubble tension is resolved at $0.006\%$ (Table~\ref{tab:hubble}): the ratio $\sqrt{\sig(120)/\sig(40)} = 1.085868$ compares to the observed ratio $1.085935$, with zero free parameter. The gravitational redshift of \PDL\ hydrogen is internally consistent with D33~\cite{PDL_D33_2026} (Proposition~\ref{prop:redshift}), and the same parameter-free architecture that fixes $\aPDL$ and the Dirac spectrum also determines $\kPDL(N)$ and $\Phicoh(r)$.

The complete \PDL\ programme is now self-consistent at five levels: the fundamental constants $\alpha$, $G$, $\mu^*$, and $H_0$ (D24, D20-V3, D28, D27); the three dynamical equations for the hydrogen atom --- Schr\"odinger (D32), Dirac (D33), and Einstein (D35); and the three falsifiable predictions $\mu^* = 1836.152670$, $(m_d - m_u)_{\PDL} = 2.532$\,MeV, and $H_{0,\mathrm{CMB}} = 67.26$\,km\,s$^{-1}$\,Mpc$^{-1}$. All emerge from the single integer quintuplet $(24, 28, 930, 10087, 11017)$.

% ================================================================
\bibliographystyle{plain}
\nocite{*}
\bibliography{References_D35}

\end{document}